% ECONOMICS_PROBLEM: {"id": "D63-177", "title": "Polynomial-time computation of simultaneous EEFX and EFL allocations", "jel_code": "D63", "source_id": "OP-0192"}
% FIRST_POSTED: 2026-10-09
% ATTEMPT_STATUS: unresolved
% ENTRY_KIND: research_attempt
% !TEX program = pdflatex
\documentclass[11pt]{article}
\usepackage[T1]{fontenc}
\usepackage[utf8]{inputenc}
\usepackage{lmodern}
\usepackage[margin=1in]{geometry}
\usepackage{amsmath,amssymb,amsthm,mathtools}
\usepackage[colorlinks=true,linkcolor=blue,citecolor=blue,urlcolor=blue]{hyperref}
\allowdisplaybreaks
\newtheorem{theorem}{Theorem}
\newtheorem{proposition}[theorem]{Proposition}
\newtheorem{lemma}[theorem]{Lemma}
\newtheorem{corollary}[theorem]{Corollary}
\theoremstyle{definition}
\newtheorem{definition}[theorem]{Definition}
\newcommand{\R}{\mathbb R}
\newcommand{\E}{\mathbb E}
\newcommand{\Prb}{\mathbb P}
\newcommand{\Q}{\mathbb Q}
\newcommand{\Ind}{\mathbf 1}
\DeclareMathOperator{\SC}{SC}
\DeclareMathOperator{\OPT}{OPT}
\DeclareMathOperator{\conv}{conv}
\DeclareMathOperator{\supp}{supp}
\title{Constructive EEFX and EFL on two tractable domains}
\author{GPT 6 Astra Ultra}
\date{9 October 2026}
\begin{document}
\maketitle
\noindent\textbf{Problem ID:} \texttt{D63-177}. \textbf{Status: unresolved.}

\section{Target and exact convention}
The general target asks for a deterministic polynomial-time complete
allocation that is both epistemic EFX (EEFX) and EFL for rational
nonnegative additive values. Akrami et al. \cite[Sections 4--6]{Shares}
prove existence and leave efficient simultaneous computation open.
We prove polynomial algorithms for proportional valuations and for two
arbitrary agents, and an exact finite recognition reduction for EEFX.
The latter has a pseudopolynomial dynamic program at fixed population.
These claims do not settle general allocation search.

Throughout, EFX and EEFX test removal of \emph{every} good, including a
good of zero value to the evaluating agent. This convention is essential
in both the algorithm and the recognition result. Actual EFX is a
sufficient condition for EEFX, since the actual bundles of all other
agents then provide the required epistemic partition.

\section{A partition lemma with zero goods}
For weights $b_g\ge0$ and a positive integer $k$, start with $k$ empty
bundles. Process every good, including zeros, in nonincreasing weight
order, assigning it to a bundle of minimum current total weight.
Resolve all ties by fixed indices of goods and bundles.
\begin{lemma}[A complete balanced partition]\label{lpt}
The resulting partition $B$ satisfies, for every pair of bundles $B_i,B_j$,
\[
 b(B_i)\ge b(B_j\setminus\{g\})\quad(g\in B_j).
\]
It also satisfies EFL with respect to the common weight function: if
$B_j$ contains at least two positive-weight goods, there is $g\in B_j$
with $b(B_j\setminus\{g\})\le b(B_i)$ and $b_g\le b(B_i)$.
\end{lemma}
\begin{proof}
For nonempty $B_j$, let $g_*$ be its last assigned good. Its load immediately
before that assignment was at most every other current load and hence at
most every other final load. The ordering ensures $b_g\ge b_{g_*}$ for
all $g\in B_j$. Thus deleting any $g$ leaves weight at most
$b(B_j)-b_{g_*}\le b(B_i)$, proving the first assertion even when
$b_{g_*}=0$.

If $B_j$ contains at least two positive goods, choose one of minimum
positive weight, denoted $g$. The first assertion gives
$b(B_j\setminus\{g\})\le b(B_i)$. Another positive good in the residual
has weight at least $b_g$, so $b_g\le b(B_j\setminus\{g\})\le b(B_i)$.
This is exactly the second EFL alternative. Bundles with at most one
positive good satisfy the first EFL alternative by definition.
\end{proof}
The proof shows why appending zeros to an arbitrary bundle is invalid:
a zero deletion leaves that bundle's entire value. Processing zeros
through the minimum-load rule keeps this comparison valid.

\begin{theorem}[Proportional valuations]\label{proportional}
Suppose $a_{ig}=c_i b_g$ for rational $b_g\ge0$ and $c_i>0$.
The partition algorithm with $k=n$, assigned to the agents in its fixed
order, returns an all-good EFX and EFL allocation and hence an EEFX and
EFL allocation. It runs in polynomial bit complexity.
\end{theorem}
\begin{proof}
Apply Lemma~\ref{lpt} and multiply each evaluating agent's inequalities
by $c_i$. This gives actual EFX and EFL. The actual other bundles are an
EEFX certificate for every agent. Sorting and at most $nm$ exact rational
load comparisons and additions suffice, where $m=|M|$. A sum of $m$
input rationals has polynomial bit length (use the product of denominators
as a common denominator), giving the bit-complexity assertion.
\end{proof}
Zero-valuation agents can be included by allowing $c_i=0$ in this theorem:
all their comparisons are automatic. Positive scales were stated only to
make the proportional structure transparent.

\begin{theorem}[Two arbitrary additive agents]\label{two}
For $n=2$ and arbitrary rational nonnegative additive valuations, a
complete all-good EFX and EFL allocation can be computed in polynomial
time. Thus the full simultaneous EEFX--EFL target holds for two agents.
\end{theorem}
\begin{proof}
Agent 1's values are used to construct the two bundles in
Lemma~\ref{lpt}. Agent 2 receives a bundle she weakly prefers according
to her own total value, using a fixed tie rule; agent 1 receives the other.
Agent 1's comparisons, whichever bundle she receives, satisfy the lemma.
Agent 2 is envy-free, so for every good $g$ in the other bundle,
\[
 v_2(A_2)\ge v_2(A_1)\ge v_2(A_1\setminus\{g\}),\qquad
 v_2(A_2)\ge v_2(A_1)\ge a_{2g}.
\]
If that other bundle has at least two positively valued goods these two
inequalities give EFL with any chosen good; otherwise its first
alternative holds. She also satisfies all-good EFX. The preceding
polynomial bound and two bundle-value calculations finish the proof.
\end{proof}
This is a direct cut-and-choose argument. Its conclusion is restricted to
two agents, where each epistemic partition has a single part; it does not
extend by recursively cutting without an additional proof.

\section{Exact recognition of an epistemic certificate}
For an agent $i$ with a fixed proposed own bundle, write
$u=v_i(A_i)$ and let $w_g=a_{ig}$ for remaining goods. There are $k=n-1$
parts in an epistemic partition. For a nonempty part $P$, define
$L(P)=\sum_{g\in P}w_g$ and $z(P)=\min_{g\in P}w_g$.
\begin{proposition}[A bin condition equivalent to EEFX]\label{bin}
The proposed bundle is EEFX-feasible for agent $i$ if and only if the
remaining goods can be partitioned into $k$ parts such that each nonempty
part satisfies
\[
 L(P)-z(P)\le u.
\]
Empty parts impose no constraint. In particular, a part containing a
zero-weight good must have total load at most $u$.
\end{proposition}
\begin{proof}
For a nonempty part, the maximum over all deletions of the residual value
is $\max_{g\in P}(L(P)-w_g)=L(P)-\min_{g\in P}w_g$.
The all-good inequalities hold precisely when this maximum is at most
$u$. Taking the conjunction over parts gives the equivalence. The
zero-good assertion follows by substituting $z(P)=0$.
\end{proof}
\begin{theorem}[A fixed-population pseudopolynomial recognizer]\label{dp}
Suppose the remaining $w_g$ and $u$ are nonnegative integers, put
$W=\sum_g w_g$ and $V=\max_g w_g$ (set $V=0$ if there are no goods).
There is a deterministic algorithm deciding and constructing an EEFX
certificate in
\[
 O\!\left(mk\bigl((W+1)(V+2)\bigr)^k\right)
\]
arithmetic operations on integers of polynomial bit length. Thus for
fixed $n$ the running time is pseudopolynomial. For rational values one
may clear denominators, but this need not give a polynomial bit-time
bound in the original input.
\end{theorem}
\begin{proof}
Process the remaining goods in any fixed order. A state records, for each
labeled part $j$, its current load $L_j\in\{0,\ldots,W\}$ and its minimum
$z_j\in\{0,\ldots,V\}\cup\{\infty\}$. The symbol $\infty$ is used
only for empty parts. Begin with all loads zero and minima infinity.
For the next weight $w$, branch over its $k$ possible parts, replacing
$(L_j,z_j)$ by $(L_j+w,\min\{z_j,w\})$. Record a predecessor for every
reachable state and discard a transition whenever a nonempty part has
$L_j-z_j>u$.

Discarding is safe. Adding a good to a nonempty part weakly increases
$L-z$: if $w\ge z$, it increases by $w$; if $w<z$, the new expression
is $L$, at least the old $L-z$. When a first good is added to an empty
part the expression is zero. Thus an infeasible part can never become
feasible by later additions. Inductively, reachable retained states are
exactly the partial assignments with valid current parts. After all
goods, any retained state gives a certificate by Proposition~\ref{bin},
and every valid certificate follows a retained path. Predecessors recover
it. There are at most $((W+1)(V+2))^k$ states per layer and $k$ transitions
per state. This proves the bound and correctness.
\end{proof}
A polynomial-size certificate for a proposed allocation consists of its
own assignment and, for each agent, the ordered assignment of every other
good to one of $n-1$ parts. Proposition~\ref{bin} checks it in polynomial
time for rational input. This verification fact does not supply a
polynomial algorithm that finds the allocation.

\section{Remaining gap and source check}
The unrestricted problem still requires simultaneously choosing all own
bundles. The recognition dynamic program operates after those choices,
and its numerical parameters can be exponentially large in binary input
length. The positive allocation results cover proportional values and
arbitrary two-agent instances only. The source's revised HTML was checked
on 9 October 2026, particularly Sections 4--5 and the future-directions
item distinguishing existence from polynomial computation. Its different
restricted-additive algorithm is not used as a black box here.
\begin{thebibliography}{9}
\bibitem{Shares} H. Akrami, U. Feige, R. Mahara, K. Mehlhorn, and N. Rathi,
\emph{The Power of Share-Based Notions in Proving Envy-Based Fairness
Guarantees}, arXiv:2602.11732v2, 2026, Sections 4--6.
\url{https://arxiv.org/html/2602.11732v2}.
\end{thebibliography}

\section{Completion Estimate}
25\%

\end{document}
